\section[Black holes and wormhole-like channels]{Black holes and wormhole-like channels: rigidity beyond the Standard Model (interface pointer)}
\label{app:bh_wormholes_pointer}

This paper focuses on the minimal stable sector and the Standard Model interface at $(m,n)=(6,3)$.
However, the same HPA--$\Omega$ rigidity philosophy extends naturally to gravitational and strong-field questions once one adopts the overhead/lapse dictionary (Section~\ref{app:time_mass_delay}) and treats boundaries as readout screens.
To keep the present manuscript self-contained, we also record (i) a minimal overhead-to-gravity closure and (ii) an executable $\chi(x)$ reconstruction protocol in Sections~\ref{app:overhead_to_gravity_closure} and \ref{app:chi_reconstruction_protocol}.
For an extended treatment of a dynamical gravity interface in the same programmatic language (routing overhead and deterministic closure), see the companion CAP-II manuscript \cite{Ma2025CAPII}.
We record here, as a pointer, two established rigidity targets.
For a deterministic finite-family calibration that links black-hole boundary information to the protocol resolution parameters $(m,n)$ used in this paper, see Appendix~\ref{app:bh_planck_capacity_calibration}.
In particular, the calibration table there provides an explicit discrete mapping $M\mapsto(m^\ast(M),n^\ast(M))$ within a bounded candidate family, which is the paper's operational notion of a black-hole ``resolution level'' at the matching layer.

\paragraph{Audit boundary (what this appendix is and is not).}
\AuditTag This appendix is a pointer module and is not used as a premise for any theorem-level folding statement in the paper.
All geometric and thermodynamic formulas quoted here (area law, Schwarzschild exterior, entropy bounds, and topological censorship) are treated as standard external inputs.
For compact assumption-level bookkeeping within this paper, see Appendix~\ref{app:physics_consensus_inputs},
Assumptions~\ref{ass:bh_planck_units_info}--\ref{ass:bh_schwarzschild_package}.
The only paper-internal contribution is the explicit protocol vocabulary used to \emph{index} these targets (boundary channel counting; overhead/lapse/delay dictionaries).
\AuditTag For the protocol-horizon / tick-trap viewpoint that separates observer-relative information horizons from GR event horizons, see Appendix~\ref{app:protocol_horizon_tick_trap} and Appendix~\ref{app:leakage_kernel}.

\paragraph{Main-text unification pointer.}
\AuditTag For the manuscript-wide unified meaning of ``wormhole-like channel'' (paid pointer-jump shortcut + four hard gates, including the local interface factorization gate LF1) and the associated audit-only operator/pole-barrier vocabulary,
see Section~\ref{subsec:x3a_wormhole_rigidity_unified_constraints}.
For the corresponding audit-facing ``black-hole-like metabolism'' hypothesis (as a coarse-grained overlay, not a theorem-level claim),
see Section~\ref{subsec:x3b_black_hole_metabolism_audit_hypothesis}.

\subsection{Black-hole area law as boundary channel counting (external input, rigid interface)}
\label{app:bh_area_law_pointer}

Black-hole thermodynamics provides a well-established link between boundary geometry and entropy: semiclassically,
\[
S_{\mathrm{BH}}=\frac{k_B A}{4\ell_P^2},
\]
where $A$ is horizon area and $\ell_P^2=G\hbar/c^3$ \cite{Bekenstein1973,BardeenCarterHawking1973,Hawking1972BlackHoles,Hawking1975,Wald1993Noether,Bousso2002}.
This is treated as a standard external input here (Assumption~\ref{ass:bh_area_law}).
This is consistent with the holographic principle viewpoint that gravitational degrees of freedom admit an effective boundary description \cite{tHooft1993,Susskind1995,Bousso2002}.
In the HPA interface language, a boundary screen is a finite-resolution readout cut with a maximal outcome count $\mathcal{N}_\partial(A,r)$, and the channel-count entropy is $S=k_B\log\mathcal{N}_\partial$ \cite{Shannon1948,CoverThomas2006}.
Under the covariant entropy bound and saturation, the area law is the saturation of boundary channel capacity \cite{Bousso2002}.
This viewpoint is consistent with the present paper's emphasis on finite alphabets and stable-sector compression: horizons are extreme instances of ``boundary-stable'' readout where channel counting dominates.

\paragraph{A $\chi$-cloud realization (paper-internal pointer).}
\InterfaceTag Within this paper, a directly drawable realization of an observer-relative horizon is obtained from the reconstructed overhead proxy field $\chi(x)$:
the budget-triggered threshold $\chi_\star$ and the induced boundary $\partial\mathcal R_\star$ define a $\chi$-horizon (Section~\ref{subsec:chi_budget_horizon_area_law}).
\MatchTag Under the saturation calibration of Assumption~\ref{ass:chi_channel_area_calibration}, the associated cloud capacity bits induce an effective area representative $A_\chi$, making the channel-count viewpoint explicit in protocol variables.
For the corresponding deterministic capacity-only table view of $I_\chi$ and the induced occupancy fraction $|{\mathcal R}_\star|/4^n$, see Table~\ref{tab:chi_horizon_budget_occupancy} (Appendix~\ref{app:generated_tables}).

\subsection{Einstein--Rosen throat and inversion continuation (wormhole-like channel)}
\label{app:wormhole_throat_pointer}

In isotropic coordinates, the Schwarzschild exterior admits an inversion symmetry and a minimal-surface throat on the time-symmetric slice, giving the classical Einstein--Rosen bridge template \cite{EinsteinRosen1935,Kruskal1960,Wald1984}.
For a minimal explicit formula package (coordinate map, metric form, and inversion), see Proposition~\ref{prop:er_isotropic_inversion}.
Conceptually, throat/bridge geometries are aligned with modern organizing principles that relate entanglement, horizons, and wormhole geometries (e.g.\ ER=EPR in appropriate settings) \cite{SusskindMaldacena2013}.

\subsection{Why these extensions are ``forced'' once rigidity is assumed (interface logic)}
\label{app:bh_wormholes_rigidity_logic}

The purpose of this appendix is not to import a full gravitational derivation into the present paper, but to record the sense in which certain gravitational structures become the standard rigid templates (and, in limited senses, unavoidable) once one adopts the same rigidity discipline used throughout this manuscript:
\begin{itemize}
\item \textbf{Boundary entropy is channel capacity.}
If horizons are treated as readout screens and entropy is treated as $\log$ of a finite channel count (channel-count entropy), then the covariant entropy bound supplies a canonical capacity bound, and saturation yields the area law coefficient $1/4$ as the rigid leading term (cf.\ \cite{Bousso2002}).
\item \textbf{Exterior geometry is unique under standard symmetry assumptions.}
In spherically symmetric vacuum regions, the Schwarzschild exterior is the standard unique template up to diffeomorphism (Birkhoff-type uniqueness; see, e.g., \cite{Birkhoff1923,Wald1984}), so any rigid ``black-hole sector'' that aims to reproduce classical tests has essentially no freedom in the exterior once the mass parameter is fixed.
\item \textbf{Endpoint avoidance naturally uses the Einstein--Rosen/Kruskal throat template.}
The maximal analytic extension of Schwarzschild contains an Einstein--Rosen bridge on a time-symmetric slice \cite{EinsteinRosen1935,Kruskal1960,Wald1984}.
If one seeks a minimal continuation that removes coordinate endpoint pathology while preserving the exterior, then using the throat/inversion structure as a gluing/continuation template is the most economical geometric move.
\item \textbf{Wormhole-like channels are topological shortcuts in readout geometry.}
In a protocol where boundary transport can be impedance-limited (delay) while bulk access can provide chord-like shortcuts, ``wormhole-like'' behavior can be interpreted as a controlled bulk-to-boundary interaction channel rather than as a violation of locality.
This viewpoint is compatible with standard wormhole terminology in GR \cite{MorrisThorne1988,Visser1995}.
\end{itemize}

\subsection{Interface closure statements (audit form)}
\label{app:bh_wormholes_interface_closure}

For completeness, we record the preceding logic in a compact ``input $\Rightarrow$ output'' form, separating standard external inputs from interface identifications.
In addition to entropy bounds, we also use standard kinematic delay/lapse templates as external targets at the matching layer: Section~\ref{app:time_mass_delay} records the operational Wigner--Smith proxy and the GR reference formulas \eqref{eq:schwarzschild_lapse} and \eqref{eq:shapiro_delay}.
Within this paper, the protocol-level overhead-to-gravity bridge that interprets these templates is recorded explicitly in Section~\ref{app:overhead_to_gravity_closure}, and an executable data protocol to reconstruct $\chi(x)$ is recorded in Section~\ref{app:chi_reconstruction_protocol}.

\paragraph{Observable hooks (minimal list).}
\MatchTag The present appendix keeps all strong-field statements at the pointer/interface level, but it does expose concrete matchable observables already audited elsewhere in the paper:
\begin{itemize}
\item \textbf{Weak-field gravitational time delay.} Shapiro-type delay and related lensing time-delay channels enter through the delay/lapse dictionary of Appendix~\ref{app:time_mass_delay} and the K4 delay registry audit (Appendix~\ref{app:k4_delay_audit}).
\item \textbf{Phase-derived dwell times.} Scattering phase shifts and WS-type delay proxies enter through the K4 scattering phase audit family (Appendices~\ref{app:k4_scattering_phase_delay_audit} and \ref{app:k4_ws_linewidth_audit}).
\end{itemize}

\paragraph{Unit boundary (avoid category errors).}
\AuditTag Capacity targets such as $I_{\mathrm{BH}}(A)$ and $I_{\mathrm{prot}}(m,n)$ are \emph{bit counts}, while delay observables $\tau(\omega)$ and lapse/redshift proxies are \emph{time} (or time-like) quantities in declared physical units.
Any relation between them in this paper is mediated only by explicit matching dictionaries (e.g.\ overhead $\kappa$ and $\chi=\log(\kappa/\kappa_0)$ in Appendix~\ref{app:delay_to_lapse_dictionary}, and the saturation calibration that maps $I_\chi$ to an effective area representative $A_\chi$ in Appendix~\ref{app:bh_planck_capacity_calibration}); no theorem-level identification of ``bits = seconds'' is assumed.

\begin{proposition}[Area law as boundary channel saturation (standard input)]
Assume (i) boundary entropy is channel-count entropy $S=k_B\log\mathcal{N}_\partial$ for a screen of area $A$, (ii) the covariant entropy bound $S\le k_BA/(4\ell_P^2)$ holds (Assumption~\ref{ass:bh_covariant_entropy_bound}), and (iii) horizons saturate this bound at leading order.
Then the Bekenstein--Hawking area law holds at leading order,
\[
S_{\mathrm{BH}}=\frac{k_BA}{4\ell_P^2},
\qquad
\mathcal{N}_\partial(A)=\exp\!\left(\frac{A}{4\ell_P^2}\right).
\]
\end{proposition}
\begin{proof}
This is the covariant entropy bound together with the channel-count definition of entropy; see, e.g., \cite{Bousso2002}.
\end{proof}

\begin{remark}[Bekenstein bound route to the coefficient $1/4$ (standard)]
The Bekenstein bound states $S\le 2\pi k_B ER/(\hbar c)$ for a system of energy $E$ contained in radius $R$ \cite{Bekenstein1981} (Assumption~\ref{ass:bh_bekenstein_bound}).
For a Schwarzschild black hole, take $E=Mc^2$ and $R=R_s=2GM/c^2$, and recall $A=4\pi R_s^2$ \cite{Wald1984}. Then
\[
S\le 2\pi k_B\frac{(Mc^2)(2GM/c^2)}{\hbar c}= \frac{4\pi k_B GM^2}{\hbar c}
=k_B\frac{A}{4\ell_P^2},
\]
where $\ell_P^2=G\hbar/c^3$. Equality is attained (at leading order) for black holes.
\end{remark}

\begin{proposition}[Bekenstein bound in the \texorpdfstring{$r$}{r}-coordinate (matching form)]
\label{prop:bekenstein_bound_r_coordinate}
Let $\mu$ be a mass scale and define the resolution coordinate $r(\mu)=\log_\varphi(\mu/m_e)$ as in \eqref{eq:r_of_mu_z128} (equivalently \eqref{eq:r_as_log_compton}).
Let $\lambda_{C,e}:=\hbar/(m_ec)$ be the electron Compton wavelength.
Then the Bekenstein bound can be written as
\[
S\ \le\ 2\pi k_B\,\frac{R}{\lambda_{C,e}}\;\varphi^{\,r(\mu)}.
\]
For a Schwarzschild black hole with $\mu=M$ and $R=R_s=2GM/c^2$, this reproduces the leading area-law scaling $S\propto \varphi^{2r(M)}$ and is equivalent to the $k_BA/(4\ell_P^2)$ form.
\end{proposition}
\begin{proof}
By the Bekenstein bound \cite{Bekenstein1981} (Assumption~\ref{ass:bh_bekenstein_bound}), $S\le 2\pi k_B ER/(\hbar c)$.
With $E=\mu c^2$, this is $S\le 2\pi k_B (\mu cR/\hbar)=2\pi k_B (R/\lambda_{C}(\mu))$, where $\lambda_C(\mu)=\hbar/(\mu c)$.
Using $\mu=m_e\varphi^{r(\mu)}$ and $\lambda_C(\mu)=\lambda_{C,e}\,(m_e/\mu)$, one obtains
\[
\frac{R}{\lambda_C(\mu)}=\frac{R}{\lambda_{C,e}}\frac{\mu}{m_e}=\frac{R}{\lambda_{C,e}}\,\varphi^{r(\mu)}.
\]
For $\mu=M$ and $R=2GM/c^2$, the bound becomes $S\le 4\pi k_B GM^2/(\hbar c)=k_BA/(4\ell_P^2)$ as in the preceding remark.
\end{proof}

\begin{remark}[Schwarzschild thermodynamic scaling in the \texorpdfstring{$r$}{r}-coordinate (external)]
For a Schwarzschild black hole, $S_{\mathrm{BH}}\propto M^2$ while the Hawking temperature scales as $T_H\propto 1/M$ \cite{Hawking1975,Wald1984} (Assumption~\ref{ass:bh_schwarzschild_package}).
Therefore, in the $r$-coordinate one has the linear log-laws
\[
\log_\varphi S_{\mathrm{BH}} = 2\,r(M) + \text{const},
\qquad
\log_\varphi T_H = -\,r(M) + \text{const}.
\]
This highlights why a logarithmic mass coordinate is a canonical matching language across microphysical scales and gravitational thermodynamics.
\end{remark}

\begin{proposition}[Schwarzschild exterior as the rigid spherically symmetric vacuum template (standard)]
In a spherically symmetric vacuum region, the Lorentzian metric is locally isometric to the Schwarzschild exterior (Birkhoff-type uniqueness), hence classical weak-field tests in such a region have no additional functional freedom beyond the mass parameter.
\end{proposition}
\begin{proof}
Standard; see, e.g., \cite{Birkhoff1923,Wald1984}.
\end{proof}

\begin{proposition}[Einstein--Rosen throat and inversion symmetry in isotropic radius (standard)]
\label{prop:er_isotropic_inversion}
Let $R_s:=2GM/c^2$ denote the Schwarzschild radius of a mass $M$.
Here $r$ denotes the Schwarzschild \emph{areal radius} coordinate (not the resolution coordinate $r(\mu)=\log_\varphi(\mu/m_e)$ used elsewhere in this paper).
Introduce the isotropic radius $\rho>0$ by the standard change of variables
\begin{equation}\label{eq:schwarzschild_isotropic_radius}
r=\rho\left(1+\frac{R_s}{4\rho}\right)^2.
\end{equation}
Then the Schwarzschild exterior metric can be written in isotropic form as
\begin{equation}\label{eq:schwarzschild_isotropic_metric}
\dd s^2
=-\left(\frac{1-\frac{R_s}{4\rho}}{1+\frac{R_s}{4\rho}}\right)^2 c^2\,\dd t^2
+\left(1+\frac{R_s}{4\rho}\right)^4\left(\dd\rho^2+\rho^2\,\dd\Omega^2\right),
\end{equation}
where $\dd\Omega^2:=\dd\theta^2+\sin^2\theta\,\dd\phi^2$.
Define $\rho_h:=R_s/4$ and the inversion map
\begin{equation}\label{eq:schwarzschild_isotropic_inversion}
\mathcal{I}:\rho\mapsto \frac{\rho_h^2}{\rho}.
\end{equation}
Then $r(\rho)=r(\mathcal{I}(\rho))$, and the isotropic form \eqref{eq:schwarzschild_isotropic_metric} is invariant under $\rho\mapsto\mathcal{I}(\rho)$.
The horizon $r=R_s$ corresponds to $\rho=\rho_h$, and the isotropic chart covers the exterior region $r\ge R_s$.
In particular, the time-symmetric spatial slice admits a two-ended completion with asymptotically flat ends ($\rho\to\infty$ and $\rho\to 0$) glued at the minimal surface $\rho=\rho_h$ (the Einstein--Rosen throat), rather than a coordinate chart of the Lorentzian black-hole interior.
\end{proposition}
\begin{proof}
The coordinate transform \eqref{eq:schwarzschild_isotropic_radius} and the isotropic form \eqref{eq:schwarzschild_isotropic_metric} are standard; see, e.g., \cite{Wald1984}.
The inversion identity $r(\rho)=r(\rho_h^2/\rho)$ follows by direct substitution into \eqref{eq:schwarzschild_isotropic_radius}.
Under $\rho=\rho_h^2/\tilde{\rho}$, one has $\dd\rho^2+\rho^2\dd\Omega^2=(\rho_h^4/\tilde{\rho}^4)(\dd\tilde{\rho}^2+\tilde{\rho}^2\dd\Omega^2)$ and $\left(1+R_s/(4\rho)\right)^4=(\tilde{\rho}+\rho_h)^4/\rho_h^4$, so the spatial conformal factor in \eqref{eq:schwarzschild_isotropic_metric} is invariant.
The throat statement is the classical Einstein--Rosen bridge template on the time-symmetric slice; see \cite{EinsteinRosen1935,Kruskal1960,Wald1984}.
\end{proof}

\begin{definition}[Wormhole-like channel as a pointer jump (protocol-level)]
\label{def:wormhole_pointer_jump}
Fix a Hilbert order $n$ and a locality-preserving address map $H_n:\{0,\dots,4^n-1\}\to\{0,\dots,2^n-1\}^2$ (Section~\ref{sec:hilbert_addressing}).
A \emph{wormhole link} is a directed (or undirected) pointer
\[
a \xrightarrow{\ \mathrm{ptr}\ } b,
\qquad
a,b\in\{0,\dots,4^n-1\},
\]
interpreted as an additional readout-level shortcut channel that bypasses the nearest-neighbor scan traversal between indices $a$ and $b$ induced by the Hilbert path.
Equivalently, it explicitly relaxes the scan adjacency constraint by augmenting the protocol with a nonlocal pointer edge in index space.
\end{definition}

\begin{remark}[Operator-mother-space encoding (dictionary only)]
\label{rem:wormhole_pointer_jump_operator_encoding}
\AuditTag The pointer-jump definition above is protocol-level.
In the operator-mother-space bookkeeping layer (Appendix~\ref{app:operator_mother_space}), an added shortcut edge may be encoded abstractly as a finite-rank (hence trace-class) perturbation $\Delta$ of a trace-class source operator $F$, so that the updated candidate is $F\mapsto F+\Delta$.
This is a dictionary identification used to keep ``wormhole-like'' language strictly at the interface level; it is not a claim about traversable Lorentzian wormholes in vacuum GR.
\end{remark}

\begin{remark}[Traversability is not assumed in the present paper]
Classical traversable wormholes in Lorentzian GR require additional conditions and are typically associated with violations of standard energy conditions; see, e.g., \cite{MorrisThorne1988,Visser1995}. For explicit traversable constructions in AdS settings via controlled boundary couplings, see \cite{GaoJafferisWall2017Traversable,MaldacenaStanfordYang2017Diving}.
Accordingly, the ``wormhole-like channel'' language used here is protocol-level and refers to a controlled bulk--boundary interaction/shortcut template in readout geometry, not to a claim of a traversable Lorentzian wormhole in vacuum GR.
In particular, the only concrete object fixed in this paper is the protocol-level pointer-jump model of Definition~\ref{def:wormhole_pointer_jump}, together with the delay/overhead dictionaries used for matching (Section~\ref{app:time_mass_delay}).
\end{remark}

\subsection{Wormhole-like pointer jump audit (protocol-only)}
\label{app:wormhole_pointer_jump_audit}

\AuditTag This subsection is a protocol-only audit: it evaluates how a finite set of explicit nonlocal pointer edges reduces scan distance on the Hilbert-indexed screen at fixed $n$.
No geometric or energy-condition assumptions are used.
The purpose is to keep ``wormhole-like'' language concrete as a discrete readout shortcut model (Definition~\ref{def:wormhole_pointer_jump}).

\begin{table}[h]
\centering
\scriptsize
\setlength{\tabcolsep}{6pt}
\renewcommand{\arraystretch}{1.15}
\begin{tabular}{r r r r r r l}
\toprule
$n$ & $4^n$ & \#ptr & mean ratio & min ratio & max $\Delta$ & exemplar \\
\midrule
3 & 64 & 0 & 1.000000 & 1.000000 & 0 & 0->63 (d0-d1=0) \\
3 & 64 & 1 & 0.790248 & 0.015873 & 62 & 0->63 (d0-d1=62) \\
3 & 64 & 2 & 0.594410 & 0.015873 & 62 & 0->63 (d0-d1=62) \\
3 & 64 & 3 & 0.466062 & 0.015873 & 62 & 0->63 (d0-d1=62) \\
3 & 64 & 4 & 0.454157 & 0.015873 & 62 & 0->63 (d0-d1=62) \\
3 & 64 & 8 & 0.374892 & 0.015873 & 62 & 0->63 (d0-d1=62) \\
4 & 256 & 0 & 1.000000 & 1.000000 & 0 & 0->255 (d0-d1=0) \\
4 & 256 & 1 & 0.824191 & 0.003922 & 254 & 0->255 (d0-d1=254) \\
4 & 256 & 2 & 0.642670 & 0.003922 & 254 & 0->255 (d0-d1=254) \\
4 & 256 & 3 & 0.499594 & 0.003922 & 254 & 0->255 (d0-d1=254) \\
4 & 256 & 4 & 0.486414 & 0.003922 & 254 & 0->255 (d0-d1=254) \\
4 & 256 & 8 & 0.413362 & 0.003922 & 254 & 0->255 (d0-d1=254) \\
5 & 1024 & 0 & 1.000000 & 1.000000 & 0 & 0->1023 (d0-d1=0) \\
5 & 1024 & 1 & 0.850936 & 0.000978 & 1022 & 0->1023 (d0-d1=1022) \\
5 & 1024 & 2 & 0.467772 & 0.000978 & 1022 & 0->1023 (d0-d1=1022) \\
5 & 1024 & 3 & 0.268401 & 0.000978 & 1022 & 0->1023 (d0-d1=1022) \\
5 & 1024 & 4 & 0.233627 & 0.000978 & 1022 & 0->1023 (d0-d1=1022) \\
5 & 1024 & 8 & 0.160657 & 0.000978 & 1022 & 0->1023 (d0-d1=1022) \\
6 & 4096 & 0 & 1.000000 & 1.000000 & 0 & 0->4095 (d0-d1=0) \\
6 & 4096 & 1 & 0.869108 & 0.000244 & 4094 & 0->4095 (d0-d1=4094) \\
6 & 4096 & 2 & 0.386594 & 0.000244 & 4094 & 0->4095 (d0-d1=4094) \\
6 & 4096 & 3 & 0.144707 & 0.000244 & 4094 & 0->4095 (d0-d1=4094) \\
6 & 4096 & 4 & 0.144707 & 0.000244 & 4094 & 0->4095 (d0-d1=4094) \\
6 & 4096 & 8 & 0.085203 & 0.000244 & 4094 & 0->4095 (d0-d1=4094) \\%
\end{tabular}
\caption{Pointer-jump reduction of scan distance on the $n$-Hilbert-indexed screen. Baseline distance is the unit-cost scan path length $d_0=|i-j|$. With a finite set of pointer edges, the shortest-path distance is $d_1$, and the ratio is $d_1/d_0$. $\Delta:=d_0-d_1$ is the absolute reduction. Rows are reproduced by \texttt{scripts/exp\_wormhole\_pointer\_jump\_audit.py}.}
\end{table}

\paragraph{Pointer-jump audit summary.} \AuditTag We model the Hilbert-indexed screen as a unit-cost scan path on indices $0,1,\dots,4^n-1$ augmented by a finite set of explicit undirected pointer edges. For a fixed finite set of index pairs we report the baseline scan distance $d_0=|i-j|$ and the shortest-path distance $d_1$ in the augmented graph, summarized by the ratio $d_1/d_0$ and the maximal absolute reduction $\Delta=d_0-d_1$. This is a protocol-only audit of readout shortcut geometry (Definition~\ref{def:wormhole_pointer_jump}).
\noindent\AuditTag For $n=3$ (screen size $4^n=64$) and \#ptr=0, the audit reports ratios over 12 fixed index pairs.
\noindent\AuditTag For $n=3$ (screen size $4^n=64$) and \#ptr=1, the audit reports ratios over 12 fixed index pairs.
\noindent\AuditTag For $n=3$ (screen size $4^n=64$) and \#ptr=2, the audit reports ratios over 12 fixed index pairs.
\noindent\AuditTag For $n=3$ (screen size $4^n=64$) and \#ptr=3, the audit reports ratios over 12 fixed index pairs.
\noindent\AuditTag For $n=3$ (screen size $4^n=64$) and \#ptr=4, the audit reports ratios over 12 fixed index pairs.
\noindent\AuditTag For $n=3$ (screen size $4^n=64$) and \#ptr=8, the audit reports ratios over 12 fixed index pairs.
\noindent\AuditTag For $n=4$ (screen size $4^n=256$) and \#ptr=0, the audit reports ratios over 12 fixed index pairs.
\noindent\AuditTag For $n=4$ (screen size $4^n=256$) and \#ptr=1, the audit reports ratios over 12 fixed index pairs.
\noindent\AuditTag For $n=4$ (screen size $4^n=256$) and \#ptr=2, the audit reports ratios over 12 fixed index pairs.
\noindent\AuditTag For $n=4$ (screen size $4^n=256$) and \#ptr=3, the audit reports ratios over 12 fixed index pairs.
\noindent\AuditTag For $n=4$ (screen size $4^n=256$) and \#ptr=4, the audit reports ratios over 12 fixed index pairs.
\noindent\AuditTag For $n=4$ (screen size $4^n=256$) and \#ptr=8, the audit reports ratios over 12 fixed index pairs.
\noindent\AuditTag For $n=5$ (screen size $4^n=1024$) and \#ptr=0, the audit reports ratios over 12 fixed index pairs.
\noindent\AuditTag For $n=5$ (screen size $4^n=1024$) and \#ptr=1, the audit reports ratios over 12 fixed index pairs.
\noindent\AuditTag For $n=5$ (screen size $4^n=1024$) and \#ptr=2, the audit reports ratios over 12 fixed index pairs.
\noindent\AuditTag For $n=5$ (screen size $4^n=1024$) and \#ptr=3, the audit reports ratios over 12 fixed index pairs.
\noindent\AuditTag For $n=5$ (screen size $4^n=1024$) and \#ptr=4, the audit reports ratios over 12 fixed index pairs.
\noindent\AuditTag For $n=5$ (screen size $4^n=1024$) and \#ptr=8, the audit reports ratios over 12 fixed index pairs.
\noindent\AuditTag For $n=6$ (screen size $4^n=4096$) and \#ptr=0, the audit reports ratios over 12 fixed index pairs.
\noindent\AuditTag For $n=6$ (screen size $4^n=4096$) and \#ptr=1, the audit reports ratios over 12 fixed index pairs.
\noindent\AuditTag For $n=6$ (screen size $4^n=4096$) and \#ptr=2, the audit reports ratios over 12 fixed index pairs.
\noindent\AuditTag For $n=6$ (screen size $4^n=4096$) and \#ptr=3, the audit reports ratios over 12 fixed index pairs.
\noindent\AuditTag For $n=6$ (screen size $4^n=4096$) and \#ptr=4, the audit reports ratios over 12 fixed index pairs.
\noindent\AuditTag For $n=6$ (screen size $4^n=4096$) and \#ptr=8, the audit reports ratios over 12 fixed index pairs.

\begin{remark}[Topological censorship constraint (standard)]
Under standard global assumptions and energy conditions (e.g.\ the null energy condition), topological censorship theorems constrain macroscopic traversable wormholes connecting asymptotically flat regions in classical GR; see, e.g., \cite{FriedmanSchleichWitt1993,Galloway1999TopologicalCensorship}.
This provides an additional reason to keep the present paper's ``wormhole-like'' language explicitly at the protocol/interface level.
\end{remark}

For the present paper, the role of this pointer is conservative: it indicates that the same rigidity discipline used here (finite alphabets, auditable closures, explicit mismatch/overhead proxies) can be extended beyond the Standard Model interface to strong-field geometry, while keeping the theorem-level folding layer clean of continuum assumptions.


